\documentclass[11pt]{article}

% Packages
\usepackage[english]{babel}
\usepackage{csquotes}
\usepackage{graphicx}
\usepackage[font=small,labelfont=bf]{caption}
\usepackage{parskip}
\usepackage{geometry}
\usepackage{setspace}


\geometry{a4paper, margin=1in}
\title{Lab 3\@: Concurrency Analysis and Logging Tool}
\author{Roy Joseph Rimachi Costillas}
\date{\today}

\singlespacing

\begin{document}

\maketitle

\section{Objective}

The objective of this activity is to extend the Lab 3 concurrency scripts with a
logging layer that detects lock waits, lock timeouts and deadlocks, records the
relevant context in a dedicated \texttt{concurrency\_log} table, and provides a
short report that summarizes the collected events by timeframe and isolation
level.

\section{Implementation}

The implementation keeps the existing Lab 3 database and adds three pieces:
the log table, SQL logging functions, and the
\texttt{log\_concurrency.sh} harness. The harness starts concurrent psql
sessions against the lab database, coordinates the acquisition of row locks,
and stores one row per detected conflict.

The activity is driven entirely from \texttt{labs/lab\_3}: the following
scripts are the deliverable, and no repository-level \texttt{make} target is
required.

\begin{center}
\begin{tabular}{ll}
\hline
File & Role \\
\hline
\texttt{scripts/create\_db\_env.sh} & Start PostgreSQL and initialize the database \\
\texttt{scripts/init\_db.sql} & Create \texttt{account} and \texttt{concurrency\_log} \\
\texttt{scripts/concurrency\_logging.sql} & Logging function and lock probes \\
\texttt{scripts/log\_concurrency.sh} & Run the scenarios or the report \\
\texttt{scripts/concurrency\_report.sql} & Timeframe and isolation-level queries \\
\texttt{scripts/simulate\_deadlock.sh} & Automated deadlock entry point \\
\hline
\end{tabular}
\end{center}

\subsection{The concurrency log table}

The \texttt{concurrency\_log} table is append-only. Its required columns are
\texttt{log\_id}, \texttt{log\_timestamp}, \texttt{isolation\_level},
\texttt{conflict\_type} and \texttt{resource\_info}. The remaining columns add
the information needed to troubleshoot a conflict: the table and row involved,
the measured wait, the PostgreSQL SQLSTATE, the backend and blocking PIDs, the
originating scenario and a JSONB details field.

\begin{verbatim}
CREATE TABLE IF NOT EXISTS concurrency_log (
    log_id           BIGSERIAL PRIMARY KEY,
    log_timestamp    TIMESTAMPTZ NOT NULL DEFAULT clock_timestamp(),
    isolation_level  TEXT NOT NULL,
    conflict_type    TEXT NOT NULL,
    resource_info    TEXT NOT NULL,
    table_name       TEXT,
    row_info         TEXT,
    wait_duration_ms NUMERIC(12, 3),
    error_code       TEXT,
    backend_pid      INT,
    blocking_pid     INT,
    scenario         TEXT,
    details          JSONB,
    logged_by        NAME NOT NULL DEFAULT current_user
);
\end{verbatim}

The column mapping to the activity requirements is summarized below.

\begin{center}
\begin{tabular}{ll}
\hline
Requirement & Column \\
\hline
Timestamp & \texttt{log\_timestamp} \\
Isolation level & \texttt{isolation\_level} \\
Conflict type & \texttt{conflict\_type} \\
Resource information & \texttt{resource\_info} \\
Table involved & \texttt{table\_name} \\
Row involved & \texttt{row\_info} \\
Wait duration & \texttt{wait\_duration\_ms} \\
PostgreSQL error code & \texttt{error\_code} \\
Blocking relationship & \texttt{backend\_pid}, \texttt{blocking\_pid} \\
\hline
\end{tabular}
\end{center}

\subsection{Detection scenarios}

Three scenarios are provided. They use \texttt{SELECT ... FOR UPDATE} to
acquire real row locks without changing the account balances.

\begin{itemize}
    \item \textbf{Lock wait.} A holder transaction locks Alice's row and sleeps
    for three seconds. A second transaction tries to lock the same row. The
    \texttt{demo\_lock\_wait\_probe()} function measures the wait with
    \texttt{clock\_timestamp()} and logs a \texttt{lock\_wait} event when the
    wait exceeds 50 milliseconds.

    \item \textbf{Lock timeout.} The holder locks the row, while the waiter sets
    a transaction-local \texttt{lock\_timeout} of one second. The
    \texttt{demo\_lock\_timeout\_probe()} function catches SQLSTATE
    \texttt{55P03} and logs a \texttt{lock\_timeout} event without aborting the
    logging transaction.

    \item \textbf{Deadlock.} Session A locks Alice and then Bob, while session B
    locks Bob and then Alice. PostgreSQL detects the cycle and aborts one
    transaction with SQLSTATE \texttt{40P01}. The harness identifies the victim
    from the psql error output and logs the victim and blocker PIDs.
\end{itemize}

Readiness does not depend on fixed delays. Each background session sets a
unique \texttt{application\_\allowbreak name}, and the harness polls
\texttt{pg\_stat\_\allowbreak activity} until that session is sleeping in
\texttt{pg\_sleep}. This confirms that the first lock has been acquired before
the second transaction starts.

\section{Results}

Figure~\ref{fig:harness} shows one complete run of the harness. The run produced
one event of each type: a lock wait of approximately 2.9 seconds, a lock
timeout of approximately 1 second with SQLSTATE \texttt{55P03}, and a deadlock
with SQLSTATE \texttt{40P01}. The final table confirms that the events were
stored in \texttt{concurrency\_log}.

\begin{figure}[htbp]
    \centering
    \includegraphics[width=\linewidth]{../images/log_concurrency.png}
    \caption{Complete \texttt{log\_concurrency.sh} run with all three scenarios
    and the latest log rows.}
    \label{fig:harness}
\end{figure}

Figure~\ref{fig:lockwait} shows the lock wait scenario in isolation. The waiter
blocked for 2927.448 milliseconds while the holder kept Alice's row locked, and
the resulting row records the isolation level, resource, wait duration and
backend PIDs.

\begin{figure}[htbp]
    \centering
    \includegraphics[width=\linewidth]{../images/lock_wait.v1.png}
    \caption{Lock wait scenario under \texttt{READ COMMITTED}.}
    \label{fig:lockwait}
\end{figure}

Figure~\ref{fig:locktimeout} shows the timeout path. PostgreSQL returned
SQLSTATE \texttt{55P03} after 1001.850 milliseconds, and the exception handler
logged the event as \texttt{lock\_timeout} instead of letting the transaction
abort without a record.

\begin{figure}[htbp]
    \centering
    \includegraphics[width=\linewidth]{../images/lock_timeout.png}
    \caption{Lock timeout scenario with SQLSTATE \texttt{55P03}.}
    \label{fig:locktimeout}
\end{figure}

Figure~\ref{fig:deadlock} shows the deadlock path. PostgreSQL selected one
session as the victim, reported \texttt{deadlock detected}, and the harness
stored the event with the victim and blocker backend PIDs and SQLSTATE
\texttt{40P01}.

\begin{figure}[htbp]
    \centering
    \includegraphics[width=\linewidth]{../images/deadlock.png}
    \caption{Deadlock scenario after PostgreSQL breaks the lock cycle.}
    \label{fig:deadlock}
\end{figure}

\section{Reporting}

The report is generated with a configurable timeframe:

\begin{verbatim}
./labs/lab_3/scripts/log_concurrency.sh --report --minutes 60
\end{verbatim}

\texttt{scripts/concurrency\_report.sql} produces four result sets: events
grouped by conflict type and isolation level, totals by conflict type, lock
hotspots by table, and the most recent events. Figure~\ref{fig:report1} shows
the first run, where each conflict type has exactly one event.

\begin{figure}[htbp]
    \centering
    \includegraphics[width=\linewidth]{../images/concurrency_report.v1.png}
    \caption{Initial concurrency report for the last 60 minutes.}
    \label{fig:report1}
\end{figure}

Figure~\ref{fig:report2} shows the report after additional scenario runs. The
counts have increased, the averages are visible, and the recent-events section
contains the complete history of the narrower timeframe.

\begin{figure}[htbp]
    \centering
    \includegraphics[width=\linewidth]{../images/concurrency_report.v2.png}
    \caption{Report after additional scenario runs.}
    \label{fig:report2}
\end{figure}

The same summary can be expressed directly in SQL:

\begin{verbatim}
SELECT conflict_type,
       isolation_level,
       count(*) AS events,
       round(avg(wait_duration_ms), 1) AS avg_wait_ms,
       round(max(wait_duration_ms), 1) AS max_wait_ms
FROM concurrency_log
WHERE log_timestamp >= now() - make_interval(mins => 60)
GROUP BY conflict_type, isolation_level
ORDER BY events DESC;
\end{verbatim}

\section{Isolation-Level Comparison}

The harness accepts an isolation level, so the same lock-wait scenario can be
repeated under \texttt{REPEATABLE READ}. Figure~\ref{fig:repeatable} shows the
result: the wait is measured and logged with isolation level
\texttt{REPEATABLE READ}. This demonstrates why storing and grouping by
\texttt{isolation\_level} is useful when comparing conflict rates across
transaction configurations.

\begin{figure}[htbp]
    \centering
    \includegraphics[width=\linewidth]{../images/lock_wait.v2.png}
    \caption{Lock wait scenario under \texttt{REPEATABLE READ}.}
    \label{fig:repeatable}
\end{figure}

\section{Production Value}

The activity demonstrates a pattern that scales beyond the lab:

\begin{itemize}
    \item Concurrency conflicts are recorded with enough context to identify
    the table, row, wait duration, isolation level and blocking session.
    \item The timeframe report supports troubleshooting and can feed a
    dashboard or an alert when deadlocks or long waits exceed a threshold.
    \item Isolation-level grouping helps decide whether a stricter isolation
    level is creating more conflicts than the application can tolerate.
    \item The append-only history supports trend analysis and database tuning;
    a production deployment would add retention or partition the table by
    timestamp.
\end{itemize}

The lab scenarios are synthetic and the harness is not a continuous background
collector. A production implementation would poll
\texttt{pg\_stat\_activity} and \texttt{pg\_locks}, or consume PostgreSQL
server logs, and would keep the monitoring load separate from application
transactions.

\section{Conclusion}

The implementation satisfies the activity requirements: it detects lock waits,
lock timeouts and deadlocks; stores timestamp, isolation level, conflict type
and resource information in \texttt{concurrency\_log}; and provides a short,
parameterized reporting query that summarizes events by timeframe and
isolation level. The stored information is sufficient to explain what waited,
on which row, for how long and which session was blocking it.

\section{Reproduction}

The complete sequence is run from the repository root and uses only the scripts
in \texttt{labs/\allowbreak lab\_3/\allowbreak scripts}:

\begin{verbatim}
./labs/lab_3/scripts/create_db_env.sh
./labs/lab_3/scripts/log_concurrency.sh
./labs/lab_3/scripts/log_concurrency.sh --report --minutes 60
\end{verbatim}

The first command starts PostgreSQL, creates \texttt{concurrency\_db}, loads
\texttt{init\_db.sql} and installs the logging functions from
\texttt{concurrency\_logging.sql}. The second command runs the three scenarios,
and the third prints the report.

Individual scenarios can be repeated with:

\begin{verbatim}
./labs/lab_3/scripts/log_concurrency.sh --scenario lock-wait
./labs/lab_3/scripts/log_concurrency.sh --scenario lock-timeout
./labs/lab_3/scripts/log_concurrency.sh --scenario deadlock
\end{verbatim}

The isolation level is selected with the
\texttt{--isolation} flag, for example
\texttt{REPEATABLE READ}.

\end{document}
